\chapter{Background}\label{chap:background}
This chapter summarizes key concepts relevant to this thesis.

\section{Wildfire Detection in Imagery}
Approaches range from handcrafted features (e.g., color-based rules) to data-driven classifiers and deep neural networks. Edge deployment motivates compact models and efficient inference.

\section{Multilayer Perceptrons}
An MLP comprises fully connected layers with nonlinear activations. For per-pixel classification, a small network can map RGB inputs to a binary fire/non-fire probability. Quantization and activation approximations (e.g., lookup tables) enable efficient hardware mapping.

\section{Fixed-Point Arithmetic and LUT Activations}
Fixed-point formats reduce memory usage and DSP requirements at modest accuracy cost. Sigmoid or similar activations can be approximated using lookup tables (LUTs) generated offline and stored in block RAM.

\section{FPGAs and VHDL}
FPGAs provide reconfigurable logic, DSP slices, and block RAMs. VHDL is used to describe the hardware. Vendor tools like Xilinx Vivado perform synthesis, placement, routing, and bitstream generation.

\section{Related Work}
Prior work has demonstrated neural inference on FPGAs using fixed-point quantization, pipelining, and memory tiling. This thesis adopts similar principles for a simple MLP.
